
\maketitle

\begin{abstract}
When does observationally coherent data determine a genuine probability
measure?  We work with a directed system of Boolean algebras carrying a
compatible family of finitely additive charges.  Coherence and finite
additivity are purely finitary conditions; the question is what further
condition licenses the passage to $\sigma$-additivity.

Finitary coherence alone does not suffice: the finite--cofinite charge
on $\mathbb{N}$ satisfies every structural condition yet assigns
persistent mass to a sequence of events that vanishes.  More generally,
$\sigma$-additivity is not first-order axiomatizable among finitely
additive probability Boolean algebras~\cite{MahonCE2026}.  Some
condition beyond finitary coherence is therefore genuinely required.

We call this condition \emph{collective exhaustion}~(CE): whenever a
decreasing sequence of events vanishes globally, there exists a level at
which the mass converges to zero.  A compatible charge family extends to
a unique $\sigma$-additive probability measure on the observable
$\sigma$-algebra if and only if it satisfies CE.

Two constructions realise the extension.  A direct Carath\'{e}odory
construction assumes compatible $\sigma$-additive marginals and builds
the global measure without topological assumptions.  A Stone
construction starts from finite additivity alone: compactness of the
Stone space supplies $\sigma$-additivity on the compact completion, and
CE reappears as the support condition ensuring the measure descends to
the realised instantiation space.  The constructions differ by where the
missing $\sigma$-additive structure is supplied --- internally by
countable reach, or externally by completion --- and whenever both
apply, they yield the same measure.

\end{abstract}

\begin{quote}\small
\textbf{2020 Mathematics Subject Classification.}
60A10 (primary); 28A60, 06E15, 28A12 (secondary).

\textbf{Keywords.}
collective exhaustion, Carathéodory observational extension, Stone observational extension, Boolean algebra, finitely additive measure, $\sigma$-additivity, directed system, Stone duality, projective limit, non-derivability, ultraproduct.
\end{quote}

% ------------------------------------------------------------------
\section{Introduction}
\label{I:sec:intro}
% ------------------------------------------------------------------

A probability measure on a field $\mathcal{F}\subseteq 2^\Omega$ is a
set function $P : \mathcal{F} \to [0,1]$
satisfying~\cite{Billingsley, Kolmogorov}:
\begin{alignat}{2}
  &\text{(positivity)}    &\qquad& P(A) \geq 0;
    \label{I:eq:positivity} \\
  &\text{(normalisation)} &\qquad& P(\Omega) = 1;
    \label{I:eq:normalisation} \\
  &\text{(countable additivity)} &\qquad&
    P\!\Bigl(\textstyle\bigcup_{n=1}^{\infty} A_n\Bigr)
    = \textstyle\sum_{n=1}^{\infty} P(A_n)
    \quad (\text{disjoint } A_n,\;
           \textstyle\bigcup A_n \in \mathcal{F}).
    \label{I:eq:countable-additivity}
\end{alignat}
For a finitely additive $P$,
\eqref{I:eq:countable-additivity} is equivalent to
\begin{equation}\label{I:eq:continuity-at-empty}
  E_n \downarrow \varnothing
  \quad\Longrightarrow\quad
  P(E_n) \to 0:
\end{equation}
mass cannot persist along a vanishing sequence of events.  In the
Kolmogorov and Carath\'{e}odory extension theorems this is
postulated~\cite{Kolmogorov, choksi1958}.

The classical extension problem already has directed structure:
compatible finite-dimensional distributions live on a directed family of
coordinate projections.  We isolate that structure --- a directed
system $(\iota, \leq)$ of event algebras $\mathcal{E}_i$ with surjective
refinement maps
\[
  \pi_{ij} : O_j \to O_i
  \qquad (i \leq j),
  \qquad
  \pi_{ik} = \pi_{ij} \circ \pi_{jk}
  \qquad (i \leq j \leq k)
\]
--- and ask what the directed structure alone determines.
Positivity and normalisation become local: each $\ell_i$ is a normalised
finitely additive charge on $\mathcal{E}_i$.  Compatibility replaces
agreement of marginals:
\begin{equation}\label{I:eq:compatibility}
  \ell_j\!\bigl(\pi_{ij}^{-1}(A)\bigr) = \ell_i(A)
  \qquad (i \leq j,\; A \in \mathcal{E}_i).
\end{equation}
These are finitary.
Countable additivity~\eqref{I:eq:countable-additivity} is not.

The finite--cofinite algebra on $\mathbb{N}$ shows how it fails.  Set
$\mu(A) = 0$ for finite $A$, $\mu(A) = 1$ for cofinite $A$.  Then
$\mu$ satisfies~\eqref{I:eq:positivity}--\eqref{I:eq:compatibility},
yet
\[
  E_n = \{n, n{+}1, \ldots\}
  \;\downarrow\; \varnothing,
  \qquad
  \mu(E_n) = 1 \;\;\forall\, n:
\]
mass persists because $\bigcap_n E_n = \varnothing$ is a countable fact
no finite check can see.  In general, $\sigma$-additivity is not
first-order axiomatizable among finitely additive probability Boolean
algebras~\cite{MahonCE2026, Hoover1978, FaginHalpernMegiddo1990}.
If $\sigma$-additivity is to be \emph{recovered} rather than assumed,
the condition must fall on the charges themselves.

We call this condition \emph{collective exhaustion}~(CE):
\begin{equation}\label{I:eq:ce}
  E_n \in \mathcal{E}_i,\quad
  E_n \downarrow \varnothing
  \quad\Longrightarrow\quad
  \exists\, j \geq i \;\text{s.t.}\;
  \ell_j\!\bigl(\pi_{ij}^{-1}(E_n)\bigr) \to 0.
\end{equation}
A compatible family extends to a unique $\sigma$-additive measure on the
observable $\sigma$-algebra if and only if it satisfies~\eqref{I:eq:ce}
(Theorem~\ref{I:thm:sp1}).  The Yosida--Hewitt
decomposition~\cite{YosidaHewitt1952} makes the content precise: $\mu =
\mu_c + \mu_p$ with $\mu_c$ countably additive and $\mu_p$ purely
finitely additive; CE is $\mu_p = 0$.

Two constructions realise the extension.  A Carath\'{e}odory
construction (\S\ref{I:sec:cara}) assumes compatible $\sigma$-additive
marginals and extends directly.  A Stone construction
(\S\ref{I:sec:stone}) starts from finite additivity alone: compactness
of the Stone space supplies $\sigma$-additivity on the completion, and
CE reappears as the support condition returning the measure to the
realised instantiation space.  Whenever both apply, they yield the same
measure (Proposition~\ref{I:prop:coincidence}).

% ------------------------------------------------------------------
\section{Observational structure and compatible charges}
\label{I:sec:setup}
% ------------------------------------------------------------------

\begin{definition}[Observational structure]
\label{I:def:obs-structure}
An \emph{observational structure} consists of:
\begin{itemize}[noitemsep]
  \item a nonempty partially ordered set $(\iota, \leq)$, the \emph{index set};
  \item for each $i \in \iota$, a measurable space $(\Out{i}, \mathcal{O}_i)$,
    the \emph{outcome space at level~$i$};
  \item for each $i \leq j$ in $\iota$, a measurable surjection
    $\pi_{ij} : \Out{j} \to \Out{i}$, the \emph{refinement map},
\end{itemize}
satisfying $\pi_{ik} = \pi_{ij} \circ \pi_{jk}$ for all $i \leq j \leq k$.
\end{definition}

\begin{definition}[Instantiation; query system]
\label{I:def:query-system}
An \emph{instantiation} is a set $\Omega$ with \emph{evaluation maps}
$\mathrm{eval}_i : \Omega \to \Out{i}$ satisfying
\[
  \pi_{ij} \circ \mathrm{eval}_j = \mathrm{eval}_i
  \qquad (i \leq j).
\]
A \emph{query system} is an observational structure together with an
instantiation.
\end{definition}

Every observational structure admits a canonical instantiation: the projective limit
\[
  \Omega
  = \varprojlim_{i \in \iota} \Out{i} = \Bigl\{
      \omega \in \prod_{i \in \iota} \Out{i}
      \;\Big|\;
      \pi_{ij}(\omega_j) = \omega_i \text{ whenever } i \leq j
    \Bigr\},
\]
with $\mathrm{eval}_i(\omega) = \omega_i$ and  $\omega \in \Omega$ generic.

\begin{example}[Bernoulli observations]
\label{I:ex:bernoulli}
Let $\iota$ be the set of finite subsets of $\mathbb{N}$, ordered by inclusion. For $i \in \iota$, set $\mathsf{O}_i = \{0,1\}^{|i|}$ with its discrete $\sigma$-algebra $\mathcal{O}_i$, and for $i \subseteq j$ let $\pi_{ij} : \mathsf{O}_j \to \mathsf{O}_i$ be the coordinate projection. An instantiation assigns to each $\omega$ a consistent family of finite binary observations; the canonical instantiation identifies $\Omega$ with $\{0,1\}^{\mathbb{N}}$.  Cylinder sets correspond to fixing finitely many coordinates, and a compatible family $\{\nu_i\}$ is precisely the system of finite-dimensional distributions of a Bernoulli process.
\end{example}

\begin{definition}[Common coarsenings]
The index set $(\iota, \leq)$ admits \emph{common coarsenings} if every pair $i, j \in \iota$ has a lower bound: there exists $k \leq i$ and $k \leq j$.
\end{definition}

\begin{definition}[Cylinder sets]
\label{I:def:cylinder}
For $i \in \iota$ and $A \in \mathcal{O}_i$, the \emph{level-$i$ cylinder with base $A$} is
\[
  \mathrm{Cyl}(i, A)
  \;=\; \{\omega \in \Omega : \mathrm{eval}_i(\omega) \in A\}.
\]
The \emph{cylinder family} is $\mathcal{C} = \{\mathrm{Cyl}(i, A) : i \in \iota,\; A \in \mathcal{O}_i\}$, and the \emph{observable $\sigma$-algebra} is $\sigma(\mathcal{C})$.
\end{definition}

\begin{lemma}[Cylinder $\pi$-system]
\label{I:lem:pi-system}
Assume $(\iota, \leq)$ admits common coarsenings.  Then $\mathcal{C}$ is a $\pi$-system: it is closed under finite intersections.
\end{lemma}

\begin{proof}
The coherence condition gives $\mathrm{Cyl}(i, A) = \mathrm{Cyl}(j, \pi_{ij}^{-1}(A))$ for $i \leq j$: the same subset of $\Omega$ is a cylinder at any finer level. Given $\mathrm{Cyl}(i, A)$ and $\mathrm{Cyl}(j, B)$, let $k$ be a common lower bound.  Then
\[
  \mathrm{Cyl}(i, A) \cap \mathrm{Cyl}(j, B)
  = \mathrm{Cyl}(k,\; \pi_{ki}^{-1}(A) \cap \pi_{kj}^{-1}(B)),
\]
which is again a cylinder.
\end{proof}

For each $i \in \iota$, let $B_i$ denote the Boolean algebra of level-$i$ cylinders $\{\mathrm{Cyl}(i, A) : A \in \mathcal{O}_i\}$; then $\mathcal{C} = \bigcup_{i \in \iota} B_i$. For $i \leq j$, define the \emph{connecting map}
\[
  \varphi_{ij} : B_i \to B_j,
  \qquad
  \varphi_{ij}(\mathrm{Cyl}(i, A)) = \mathrm{Cyl}(j, \pi_{ij}^{-1}(A)).
\]
The family $\{B_i, \varphi_{ij}\}$ is a directed system of Boolean algebras.

\begin{definition}[Compatible family]
A \emph{charge} on a Boolean algebra $B$ is a finitely additive function $\mu : B \to [0, \infty)$ with $\mu(\emptyset) = 0$.  A family of charges $\{\mu_i\}_{i \in \iota}$ with $\mu_i$ on $B_i$ is \emph{compatible} if
\[
  \mu_j(\varphi_{ij}(E)) = \mu_i(E)
  \qquad \text{for all } i \leq j \text{ and all } E \in B_i.
\]
\end{definition}

\noindent Compatibility ensures a well-defined charge on $\mathcal{C}$
via $\mu(\mathrm{Cyl}(i, A)) := \mu_i(\mathrm{Cyl}(i, A))$.

\begin{definition}[Surjective Evaluation]
\label{I:def:eval-surjective}
The query system satisfies \emph{Surjective Evaluation} if $\mathrm{eval}_i : \Omega \to \Out{i}$ is surjective for every $i \in \iota$.
\end{definition}

\begin{lemma}[Surjectivity and injectivity from Surjective Evaluation]
\label{I:lem:surj-inj}
If the query system satisfies Surjective Evaluation, then for every $i \leq j$:
\begin{enumerate}[label=(\roman*)]
  \item the refinement map $\pi_{ij} : \Out{j} \to \Out{i}$ is surjective;
  \item the connecting map $\varphi_{ij} : B_i \to B_j$, defined by $\varphi_{ij}(\mathrm{Cyl}(i,A)) = \mathrm{Cyl}(j, \pi_{ij}^{-1}(A))$, is an injective Boolean algebra homomorphism.
\end{enumerate}
\end{lemma}

\begin{proof}
For (i): given $y \in \Out{i}$, surjectivity of $\mathrm{eval}_i$ yields $\omega \in \Omega$ with $\mathrm{eval}_i(\omega) = y$. Then $\pi_{ij}(\mathrm{eval}_j(\omega)) = \mathrm{eval}_i(\omega) = y$ by the coherence condition, so $\pi_{ij}$ is surjective.

For (ii): $\varphi_{ij}$ is a Boolean homomorphism by construction. If $\varphi_{ij}(\mathrm{Cyl}(i,A)) = \varphi_{ij}(\mathrm{Cyl}(i,A'))$, then $\pi_{ij}^{-1}(A) = \pi_{ij}^{-1}(A')$ in $\Out{j}$. Applying $\pi_{ij}$ and using surjectivity gives $A = A'$, so $\varphi_{ij}$ is injective.
\end{proof}

% ------------------------------------------------------------------
\section{Direct extension from compatible probability marginals}
\label{I:sec:cara}
% ------------------------------------------------------------------

\begin{proposition}[Observational determination]
\label{I:thm:obs-determination}
Let the query system admit common coarsenings and let $P$, $P'$ be probability measures on $(\Omega, \sigma(\mathcal{C}))$.  If $(\mathrm{eval}_i)_\# P = (\mathrm{eval}_i)_\# P'$ for all $i \in \iota$, then $P = P'$.
\end{proposition}

\begin{proof}
The equality of pushforward laws implies $P$ and $P'$ agree on every cylinder event.  By Lemma~\ref{I:lem:pi-system} the cylinder family $\mathcal{C}$ is a $\pi$-system generating $\sigma(\mathcal{C})$.  The $\pi$--$\lambda$ theorem gives $P = P'$.
\end{proof}

\begin{definition}[Common refinements]
The index set $(\iota, \leq)$ admits \emph{common refinements} if every pair $i, j \in \iota$ has an upper bound: there exists $k \geq i$ and $k \geq j$.
\end{definition}

\begin{lemma}[Finite observational content]
\label{I:thm:finite-content}
Assume the query system admits common coarsenings and common refinements and satisfies Surjective Evaluation, and let $\{\nu_i\}_{i \in \iota}$ be a family with each $\nu_i \in \mathcal{P}(\Out{i})$ compatible under the refinement maps, i.e.\ $(\pi_{ij})_\# \nu_j = \nu_i$ for all $i \leq j$. Then the assignment
\[
  \mu_0(\mathrm{Cyl}(i, A)) := \nu_i(A),
  \qquad A \in \mathcal{O}_i,
\]
defines a well-defined finitely additive content on $\mathcal{C}$.
\end{lemma}

\begin{proof}
\textit{Well-definedness.}
Suppose $\mathrm{Cyl}(i, A) = \mathrm{Cyl}(j, B)$.  Use common refinements to find $k$ with $k \geq i$ and $k \geq j$.  The constraint sets in $\Out{k}$ are $\pi_{ik}^{-1}(A)$ and $\pi_{jk}^{-1}(B)$.  The cylinder equality and surjective evaluation (surjectivity of $\mathrm{eval}_k$) imply these coincide in $\Out{k}$.  Compatibility then gives $\nu_i(A) = \nu_k(\pi_{ik}^{-1}(A)) = \nu_k(\pi_{jk}^{-1}(B)) = \nu_j(B)$.

\textit{Finite additivity.}
For disjoint cylinders at possibly different levels, compress to a common refinement $k$ and use additivity of $\nu_k$.
\end{proof}

\begin{definition}[Sequential common refinements]
The index set $(\iota, \leq)$ admits \emph{sequential common refinements} if every countable family $\{i_n\}_{n \geq 1} \subset \iota$ has an upper bound in $\iota$. Sequential common refinements imply common refinements.
\end{definition}

\begin{theorem}[Carathéodory observational extension]
\label{I:thm:obs-extension}
Suppose:
\begin{enumerate}[label=(\roman*)]
  \item $(\iota, \leq)$ admits common coarsenings and sequential common
        refinements;
  \item the query system satisfies Surjective Evaluation;
  \item each $\nu_i$ is a probability measure on $(\Out{i}, \mathcal{O}_i)$,
        and the family is compatible under the refinement maps: $(\pi_{ij})_\# \nu_j = \nu_i$ for all $i \leq j$.
\end{enumerate}
Then there exists a unique probability measure $P$ on $(\Omega, \sigma(\mathcal{C}))$ such that $(\mathrm{eval}_i)_\# P = \nu_i$ for every $i \in \iota$.
\end{theorem}

\begin{proof}
\textit{Step 1: Finite content.}
Lemma~\ref{I:thm:finite-content} gives a well-defined finitely additive content $\mu_0$ on $\mathcal{C}$.

\textit{Step 2: $\sigma$-subadditivity.}
Let $\mathrm{Cyl}(i, B) \subseteq \bigcup_{n \geq 1} \mathrm{Cyl}(i_n, A_n)$. Apply sequential common refinements to $\{i, i_1, i_2, \ldots\}$ to find a single index $k$ with $k \geq i$ and $k \geq i_n$ for all $n$.  Compress each cylinder to level $k$:
\[
  \mathrm{Cyl}(i, B) = \mathrm{Cyl}(k, E),
  \qquad
  \mathrm{Cyl}(i_n, A_n) = \mathrm{Cyl}(k, E_n),
\]
where $E = \pi_{ik}^{-1}(B)$ and $E_n = \pi_{i_n k}^{-1}(A_n)$. Surjective evaluation and the cylinder inclusion in $\Omega$ imply $E \subseteq
\bigcup_n E_n$ in $\Out{k}$.  Therefore
\[
  \mu_0(\mathrm{Cyl}(i, B))
  = \nu_k(E)
  \leq \sum_{n=1}^\infty \nu_k(E_n)
  = \sum_{n=1}^\infty \mu_0(\mathrm{Cyl}(i_n, A_n)),
\]
using $\sigma$-subadditivity of the single measure $\nu_k$.

\textit{Step 3: Carathéodory extension.}
The family $\mathcal{C}$ is a set semiring generating $\sigma(\mathcal{C})$. It is closed under finite intersections by Lemma~\ref{I:lem:pi-system}. Moreover, if $E = \mathrm{Cyl}(i, A)$ and $F = \mathrm{Cyl}(j, B)$, choose $k$ with $k \geq i, j$; then
\[
  E = \mathrm{Cyl}\!\bigl(k,\, \pi_{ik}^{-1}(A)\bigr),
  \qquad
  F = \mathrm{Cyl}\!\bigl(k,\, \pi_{jk}^{-1}(B)\bigr),
\]
so
\[
  E \setminus F = \mathrm{Cyl}\!\bigl(k,\, \pi_{ik}^{-1}(A) \setminus \pi_{jk}^{-1}(B)\bigr)
  \in \mathcal{C}.
\]
Thus $\mathcal{C}$ is closed under relative differences, confirming the semiring property. Since $\mu_0$ is a finitely additive content on $\mathcal{C}$ that is $\sigma$-subadditive for countable covers by cylinder sets (Step~2), the semiring form of Carathéodory's extension theorem yields a measure $P$ on $(\Omega, \sigma(\mathcal{C}))$ extending $\mu_0$ (Bogachev~\cite{Bogachev}, Vol.~I, Corollary~1.11.9).

\textit{Step 4: Marginal recovery and normalization.}
For any $A \in \mathcal{O}_i$, $P(\mathrm{Cyl}(i, A)) = \mu_0(\mathrm{Cyl}(i, A)) = \nu_i(A)$, so $(\mathrm{eval}_i)_\# P = \nu_i$.  Setting $A = \Out{i}$ gives $P(\Omega) = 1$.

\textit{Uniqueness.}
Two probability measures with the same evaluation marginals agree on $\mathcal{C}$, hence on $\sigma(\mathcal{C})$ by Proposition~\ref{I:thm:obs-determination}.
\end{proof}

\begin{remark}[Scope and relation to Kolmogorov's extension theorem]
\label{I:rem:kolmogorov}
Theorem~\ref{I:thm:obs-extension} requires sequential common refinements: every countable family of indices must have an upper bound.  This holds for the linearly ordered index sets $(\mathbb{N}, \leq)$ and $(\mathbb{R}, \leq)$, but \emph{not} for the finite-subset indexing of Example~\ref{I:ex:bernoulli}, where a countable family of finite subsets can have infinite union.  Classical discrete-time Bernoulli processes are covered by Kolmogorov's theorem~\cite{Kolmogorov} directly; the present theorem addresses query systems whose index structure already admits countable refinements.  The Stone construction (\S\ref{I:sec:stone}) avoids sequential common refinements by using compactness, at the cost of requiring discriminability and CE.

The key conceptual difference from Kolmogorov is that no topology is imposed on the outcome spaces: the measure $P$ is built directly on the measurable structure of $\Omega$ by Carathéodory, without any appeal to tightness or compactness.
\end{remark}

% ------------------------------------------------------------------
\section{Collective Exhaustion}
\label{I:sec:ce}
% ------------------------------------------------------------------

\begin{example}[Finite--cofinite obstruction]
\label{I:ex:finite-cofinite}
On $\mathbb{N}$, let $\mathcal{E}$ be the finite--cofinite algebra and
$\mu_0(A) = 0$ for finite $A$, $\mu_0(A) = 1$ for cofinite $A$.  Then
$\mu_0$ is a normalised finitely additive charge, but
$E_n = \{n, n{+}1, \ldots\} \downarrow \varnothing$ with
$\mu_0(E_n) = 1$ for all~$n$.
\end{example}

\begin{definition}[Collective Exhaustion]
\label{I:def:ce}
A charge $\mu$ on $\mathcal{C}$ satisfies \emph{collective exhaustion}
(CE) if $E_n \downarrow \varnothing$ in $\mathcal{C}$ implies
$\mu(E_n) \to 0$.
\end{definition}

\noindent We work at the level of an abstract directed system, writing
$\mathcal{E}_i$ for the algebra at level $i$ and $\ell_i$ for its charge.

\begin{definition}[Boolean charge space]
A \emph{Boolean charge space} is a pair $(\mathcal{E}, \ell)$ where $\mathcal{E}$ is a Boolean algebra of subsets of a set $O$ and $\ell : \mathcal{E} \to [0,1]$ is a normalized finitely-additive valuation. No $\sigma$-algebra and no topology are assumed on $O$.
\end{definition}

By Stone's representation theorem~\cite{Stone1936}, $\mathcal{E} \cong
\mathrm{Clop}(S(\mathcal{E}))$ for a compact totally disconnected
$S(\mathcal{E})$, and $\sigma(\mathcal{E})$ is its Baire
$\sigma$-algebra.

\begin{lemma}[Extension criterion]
\label{I:thm:extension-criterion}
A normalized finitely-additive $\ell : \mathcal{E} \to [0,1]$ extends
uniquely to a $\sigma$-additive measure on $\sigma(\mathcal{E})$ if and
only if
\[
  E_n \downarrow \varnothing \text{ in } \mathcal{E}
  \quad\Longrightarrow\quad
  \ell(E_n) \to 0.
\]
\end{lemma}

\begin{proof}
Standard; see Halmos~\cite{Halmos1950} or
Fremlin~\cite{Fremlin2003}, Vol.~1.
\end{proof}

\noindent Confirming $\bigcap_n E_n = \emptyset$ requires access to
$\sigma(\mathcal{E})$; the condition is not verifiable within
$\mathcal{E}$ alone.

\begin{definition}[Nontrivial refinement]
Let $\pi : O_j \to O_i$ be a surjective refinement map, inducing the Boolean algebra homomorphism $\pi^{-1} : \mathcal{E}_i \to \mathcal{E}_j$.  The refinement is \emph{nontrivial} if there exists $A \in \mathcal{E}_j$ with $A \notin \pi^{-1}(\mathcal{E}_i)$.
\end{definition}

\begin{proposition}[Nontrivial refinement introduces new events]
\label{I:prop:nontrivial}
If $j \geq i$ is a nontrivial refinement, then $\pi^{-1}(\mathcal{E}_i)
\subsetneq \mathcal{E}_j$.  Equivalently, a level $i$ at which
$\mathcal{E}_j = \pi^{-1}(\mathcal{E}_i)$ for every $j \geq i$ is one above which the refinement order is algebraically trivial: no finer level adds new discriminative power.
\end{proposition}

\begin{proof}
Nontriviality asserts the existence of $A \in \mathcal{E}_j$ with $A \notin \pi^{-1}(\mathcal{E}_i)$, directly giving the strict inclusion. The contrapositive is the stated equivalence.
\end{proof}

\begin{lemma}[Strict enlargement under $\sigma$-closure]
\label{I:prop:gap-nonempty}
Let $\mathcal{E}$ be a Boolean algebra of subsets of a countably infinite set $O$ that is not a $\sigma$-algebra.  Then $\sigma(\mathcal{E}) \setminus \mathcal{E} \neq \emptyset$.
\end{lemma}

\begin{proof}
If $\mathcal{E}$ is not a $\sigma$-algebra, there exists a countable family $\{A_n\} \subset \mathcal{E}$ with $\bigcup_n A_n \notin \mathcal{E}$.  By definition $\bigcup_n A_n \in \sigma(\mathcal{E})$, so it lies in $\sigma(\mathcal{E}) \setminus \mathcal{E}$.
\end{proof}

\begin{definition}[Collectively exhaustive family]
\label{I:def:collectively-exhaustive}
A compatible family $\{\ell_i\}$ of normalized finitely-additive charges
on a directed system is \emph{collectively exhaustive} if for every
$i \in \mathcal{I}$ and every $E_n \downarrow \varnothing$ in
$\mathcal{E}_i$,
\[
  \exists\, j \geq i \;\text{such that}\;
  \ell_j\!\bigl(\pi_{ij}^{-1}(E_n)\bigr) \to 0.
\]
\end{definition}

\begin{proposition}[Index-layer conditions do not force $\sigma$-additivity]
\label{I:prop:independence}
There exists a compatible family of normalized finitely-additive charges on a sequentially upper-directed directed system that fails $\sigma$-additive extension at every level.
\end{proposition}

\begin{proof}
We take the simplest sequentially upper-directed index set: adjoin a single terminal element $\omega$ to $\mathbb{N}$, setting $n \leq \omega$ for all $n \in \mathbb{N}$.  Take all outcome spaces to be $\mathbb{N}$, all event algebras to be the finite-cofinite algebra $\mathcal{E}$, all refinement maps to be the identity, and each $\ell_i$ to be the finite-cofinite charge $\ell$ of Example~\ref{I:ex:finite-cofinite}.  Compatibility holds trivially.  The system satisfies every structural condition.  By Example~\ref{I:ex:finite-cofinite}, $\ell$ is not $\sigma$-additive and admits no $\sigma$-additive extension; hence every level fails $\sigma$-additive extension.
\end{proof}

\begin{theorem}[CE Characterisation]
\label{I:thm:sp1}
Let $\{\ell_i\}$ be a normalized compatible family of finitely-additive charges on a directed system indexed by $\mathcal{I}$.  The following are equivalent:
\begin{enumerate}[label=(\roman*)]
  \item The family is collectively exhaustive.
  \item $\ell_i$ extends uniquely to a $\sigma$-additive measure on
        $\sigma(\mathcal{E}_i)$ for every $i \in \mathcal{I}$.
\end{enumerate}
\end{theorem}

\begin{proof}
\emph{(i) $\Rightarrow$ (ii).}
Fix $i \in \mathcal{I}$ and a decreasing sequence $E_n \searrow \emptyset$ in $\mathcal{E}_i$.  By CE, find $j \geq i$ with $\ell_j(\pi_{ij}^{-1}(E_n)) \to 0$. By compatibility, $\ell_i(E_n) = \ell_j(\pi_{ij}^{-1}(E_n)) \to 0$.  So $\ell_i$ is continuous at $\emptyset$.  Lemma~\ref{I:thm:extension-criterion} gives the unique $\sigma$-additive extension.

\emph{(ii) $\Rightarrow$ (i).}
Suppose $\ell_i$ extends to $\mu_i$ for every $i$.  Let $E_n \searrow \emptyset$ in $\mathcal{E}_i$.  Then $\ell_i(E_n) = \mu_i(E_n) \to \mu_i(\emptyset) = 0$ by $\sigma$-additivity of $\mu_i$ and normalization.  Take $j = i$.
\end{proof}

The Yosida--Hewitt decomposition~\cite{YosidaHewitt1952} writes every
bounded charge uniquely as
\[
  \ell = \ell_c + \ell_p,
  \qquad
  \ell_c \text{ countably additive},
  \quad
  \ell_p \text{ purely finitely additive}.
\]
By Theorem~\ref{I:thm:sp1}, CE is $\ell_p = 0$.

\begin{proposition}[CE non-derivability]
\label{I:prop:ce-irred}
No condition expressible in the first-order language of query systems implies collective exhaustion.  More generally, $\sigma$-additivity is not first-order axiomatizable among finitely additive probability Boolean algebras~\cite{MahonCE2026}.
\end{proposition}

\begin{proof}
Proposition~\ref{I:prop:independence} provides the particular case: the finite-cofinite charge satisfies all named structural hypotheses and fails CE.  The universal case is proved in~\cite{MahonCE2026} via \L{}o\'{s}'s theorem~\cite{Los1955}: Dirac masses $\delta_n$ on $\mathcal{P}(\mathbb{N})$ are individually $\sigma$-additive, but the charge induced on the diagonal copy of $\mathcal{P}(\mathbb{N})$ in their ultraproduct is the $\{0,1\}$-valued ultrafilter charge, which fails $\sigma$-additivity.  Since a first-order theory satisfied by every $\delta_n$ is satisfied by the ultraproduct, no such theory can axiomatize $\sigma$-additivity.
\end{proof}

\begin{example}[Propositional probability as a query system]
\label{I:ex:propositional}
The framework is not specific to empirical measurement.  Let $T$ be a consistent
countable propositional theory with sentences $\phi_1,\phi_2,\ldots$  At level
$n$, let
\[
  \mathcal{E}_n \;=\; \mathrm{Lind}(\phi_1,\ldots,\phi_n)
\]
be the finite Lindenbaum algebra of Boolean combinations of
$\phi_1,\ldots,\phi_n$, modulo $T$-equivalence.  Restriction from the first $m$
sentences to the first $n$ gives the refinement map for $n \leq m$, and the
canonical instantiation is
\[
  \Omega \;=\; \varprojlim_n\, \mathrm{atoms}(\mathcal{E}_n)
  \;=\; \{\text{complete consistent extensions of } T\},
\]
with evaluation given by restriction to finite fragments.  Surjective evaluation
is exactly the Lindenbaum extension lemma: every finite consistent assignment
extends to a complete consistent extension of $T$.

A compatible family of normalised charges is therefore a coherent system of
finite-level probability assignments.  A decreasing sequence of cylinder events
has empty intersection in $\Omega$ exactly when the corresponding finite
propositional conditions are jointly unsatisfiable over $T$.  CE requires the
assigned masses to drain away in this case,
\[
  \nu_{n_k}(\alpha_k) \;\longrightarrow\; 0,
\]
ruling out persistent mass on finite conditions that collectively vanish.
Theorem~\ref{I:thm:sp1} then yields a unique $\sigma$-additive probability
measure on complete consistent extensions exactly when CE holds.
\end{example}

% ------------------------------------------------------------------
\section{Stone realization from finitely additive data}
\label{I:sec:stone}
% ------------------------------------------------------------------

\begin{proposition}[Stone duality for the direct limit]
\label{I:prop:stepA}
The Stone space of the direct limit satisfies $\mathrm{St}\!\bigl(\bigcup_i B_i\bigr) \cong \varprojlim_i \mathrm{St}(B_i)$.
\end{proposition}

\begin{proof}
By Lemma~\ref{I:lem:surj-inj}(ii), the connecting maps $\varphi_{ij}$ are injective Boolean algebra homomorphisms.  Under Stone duality, injective Boolean algebra homomorphisms correspond to surjective continuous maps between Stone spaces. The Stone space functor $\mathrm{St}$ converts colimits in the category of Boolean algebras to limits in the category of compact Hausdorff spaces
\cite[II.4]{johnstone1982}.  Applied to $\bigcup_i B_i$ with its injective
system, this gives $\mathrm{St}\!\bigl(\bigcup_i B_i\bigr) \cong \varprojlim_i \mathrm{St}(B_i)$.
\end{proof}

\label{I:sec:invlim}

\begin{proposition}[Inverse limit measure]
\label{I:prop:invlim}
Let $\{\mu_i\}$ be a compatible family of normalised charges on $\{B_i\}$. Then there exists a unique probability measure $\hat{\mu}$ on $\mathcal{B}(\varprojlim_i \mathrm{St}(B_i))$ whose pushforward along each projection $\varprojlim_i \mathrm{St}(B_i) \to \mathrm{St}(B_j)$ agrees with the Borel measure on $\mathrm{St}(B_j)$ corresponding to $\mu_j$.
\end{proposition}

\begin{proof}
\textit{Step 1: Single-algebra correspondence.}
For each $i$, Stone duality identifies $B_i$ with the clopen algebra $\mathrm{Clop}(\mathrm{St}(B_i))$.  A charge $\mu_i$ on $B_i$ thus defines a finitely additive set function on the clopens of the compact Hausdorff space $\mathrm{St}(B_i)$.  Since $\mathrm{St}(B_i)$ is compact and totally disconnected, the clopens form a base for the topology and a finitely additive charge on them extends uniquely to a regular Borel measure; see
\cite{cardona2025}, \S6, or \cite{Halmos1950}, \S\S53--54.  Call this measure
$\hat{\mu}_i$ on $\mathrm{St}(B_i)$.

\textit{Step 2: Compatibility.}
The compatibility condition $\mu_j(\varphi_{ij}(E)) = \mu_i(E)$ translates, under Stone correspondence, to $\hat{\mu}_i$ being the pushforward of $\hat{\mu}_j$ along the bonding map $\mathrm{St}(B_j) \to \mathrm{St}(B_i)$. The family $\{\hat{\mu}_i\}$ is compatible on the inverse system $\{\mathrm{St}(B_i)\}$.

\textit{Step 3: Extension to the inverse limit.}
Each $\mathrm{St}(B_i)$ is compact Hausdorff.  The inverse limit $\varprojlim_i \mathrm{St}(B_i)$ is compact Hausdorff.  By
\cite[Theorem~1]{choksi1958}, a compatible family of regular Borel probability
measures on a cofiltered inverse system of compact Hausdorff spaces with surjective bonding maps extends uniquely to a regular Borel probability measure on the inverse limit.  This gives $\hat{\mu}$ on $\mathcal{B}(\varprojlim_i \mathrm{St}(B_i))$.
\end{proof}

\begin{definition}[Discriminability]
\label{I:def:discriminability}
The query system \emph{discriminates points} if for every $\omega \neq \omega'$ in $\Omega$ there exists $E \in \mathcal{C}$ with $\omega \in E$ and $\omega' \notin E$.
\end{definition}

\noindent Write $\mathrm{pure}(\omega) = \{E \in \mathcal{C} : \omega \in E\}$ for the principal ultrafilter at $\omega$, and $[E] = \{u \in \mathrm{St}(\mathcal{C}) : E \in u\}$ for the basic clopen.  Discriminability makes $\mathrm{pure}$ injective.

\begin{proposition}[Structural identification]
\label{I:prop:B1}
Suppose the query system discriminates points.  Then
\[
  \mathrm{pure}^{-1}\!\bigl(\mathcal{B}_{\mathrm{Ba}}(\mathrm{St}(\mathcal{C}))\bigr)
  \;=\; \sigma(\mathcal{C}),
\]
where $\mathcal{B}_{\mathrm{Ba}}$ denotes the Baire $\sigma$-algebra.
\end{proposition}

\begin{proof}
\textbf{($\supseteq$):}
For any $E = \mathrm{Cyl}(i, A) \in \mathcal{C}$, $\mathrm{pure}^{-1}([E]) = \{\omega : E \in \mathrm{pure}(\omega)\} = \{\omega : \omega \in E\} = E$. Since $[E]$ is clopen, hence Baire, every cylinder set is in $\mathrm{pure}^{-1}(\mathcal{B}_{\mathrm{Ba}})$, and therefore $\sigma(\mathcal{C}) \subseteq \mathrm{pure}^{-1}(\mathcal{B}_{\mathrm{Ba}})$.

\textbf{($\subseteq$):}
The Stone space $\mathrm{St}(\mathcal{C})$ is compact and totally disconnected, so its clopens form a topological basis.  The clopens are exactly the sets $[E]$ for $E \in \mathcal{C}$ \cite[I.3.2]{johnstone1982}, so the Baire $\sigma$-algebra of $\mathrm{St}(\mathcal{C})$ is $\sigma(\{[E] : E \in \mathcal{C}\})$. For any Baire set $V$, $\mathrm{pure}^{-1}(V) \in
\sigma(\{\mathrm{pure}^{-1}([E]) : E \in \mathcal{C}\}) = \sigma(\mathcal{C})$.

\textbf{Injectivity.}
Discriminability (Definition~\ref{I:def:discriminability}) implies $\mathrm{pure}$ is injective: if $\omega \neq \omega'$, there exists $E \in \mathcal{C}$ with $\omega \in E$ and $\omega' \notin E$, so $E \in \mathrm{pure}(\omega)$ but $E \notin \mathrm{pure}(\omega')$. Injectivity ensures the pullback does not conflate distinct events.
\end{proof}

\begin{remark}
On a compact Hausdorff space the Baire $\sigma$-algebra is contained in the Borel $\sigma$-algebra, with equality when the space is second-countable. We work throughout with the Baire $\sigma$-algebra on $\mathrm{St}(\mathcal{C})$; this suffices because the measure $\hat{\mu}$ constructed above is a Baire measure (it is the unique extension of a charge on the clopen algebra), and Choksi's theorem applies to Baire measures on compact spaces.
\end{remark}

\begin{proposition}[Support condition]
\label{I:prop:B2}
Let $\mu$ be a charge on $\mathcal{C}$ satisfying CE.  Then the corresponding Baire measure $\hat{\mu}$ on $\mathrm{St}(\mathcal{C})$ is supported on $\mathrm{pure}(\Omega)$, the set of principal ultrafilters.
\end{proposition}

\begin{proof}
By the Yosida--Hewitt theorem \cite{YosidaHewitt1952, rao1983}, every bounded charge on a Boolean algebra decomposes uniquely as $\mu = \mu_c + \mu_p$, where $\mu_c$ is countably additive and $\mu_p$ is purely finitely additive.  By Theorem~\ref{I:thm:sp1}, CE is equivalent to $\mu_p = 0$, i.e.\ $\mu$ is itself $\sigma$-additive on $\mathcal{C}$.

A $\sigma$-additive probability on a Boolean algebra $\mathcal{C}$ extends uniquely to a measure on $\sigma(\mathcal{C})$.  Under the Stone correspondence, this measure lifts to a Baire measure $\hat{\mu}$ on $\mathrm{St}(\mathcal{C})$.  For any non-principal ultrafilter $u \in \mathrm{St}(\mathcal{C}) \setminus \mathrm{pure}(\Omega)$, every clopen neighbourhood $[E]$ of $u$ satisfies $E \neq \varnothing$ in $\mathcal{C}$; but because $u$ is non-principal, for each $\omega \in \Omega$ there exists $E_\omega \in \mathcal{C}$ with $\omega \in E_\omega$ and $E_\omega \notin u$.  In particular, $u$ is not an evaluation ultrafilter.  Discriminability ensures that the principal ultrafilters separate elements of $\mathcal{C}$, so $\mathrm{pure}(\Omega)$ is dense in $\mathrm{St}(\mathcal{C})$; but the support of a $\sigma$-additive measure need not concentrate on the dense set without further argument.

The required link is classical~\cite{rao1983}, Ch.~10: for a $\sigma$-additive charge $\mu$ on a Boolean algebra $\mathcal{C}$ realised as a field of subsets of $\Omega$, the corresponding Baire measure on $\mathrm{St}(\mathcal{C})$ assigns full mass to those ultrafilters that are closed under countable intersection in $\mathcal{C}$.  Each principal ultrafilter $\mathrm{pure}(\omega)$ is closed under countable intersection (it evaluates at a point of $\Omega$), so $\hat{\mu}$ is supported on $\mathrm{pure}(\Omega)$.
\end{proof}

\begin{theorem}[Stone observational extension]
\label{I:thm:stone-main}
Let the query system satisfy Surjective Evaluation, Discriminability, and common coarsenings, and let $\{\mu_i\}$ be a compatible family of normalised charges on $\{B_i\}$ satisfying CE.  Then there exists a unique probability measure $P$ on $(\Omega, \sigma(\mathcal{C}))$ such that
\[
  P(\mathrm{Cyl}(i, A))
  \;=\; \mu_i(\mathrm{Cyl}(i, A))
  \qquad \text{for all } i \in \iota \text{ and } A \in \mathcal{O}_i.
\]
\end{theorem}

\begin{proof}
The proof assembles the preceding propositions:
\[
  \{\mu_i\}
  \;\xrightarrow{\S\ref{I:sec:invlim},\,\text{Step~1}}\;
  \{\hat{\mu}_i\}
  \;\xrightarrow{\S\ref{I:sec:invlim},\,\text{Step~3}}\;
  \hat{\mu} \text{ on } \varprojlim_i \mathrm{St}(B_i)
  \;\xrightarrow{\text{Prop.}~\ref{I:prop:B1}}\;
  P \text{ on } \sigma(\mathcal{C}).
\]

\textit{Existence.}
By Proposition~\ref{I:prop:stepA}, $\mathrm{St}(\mathcal{C}) \cong \varprojlim_i \mathrm{St}(B_i)$. By Proposition~\ref{I:prop:invlim}, the compatible family $\{\mu_i\}$ extends to a Baire probability measure $\hat{\mu}$ on $\varprojlim_i \mathrm{St}(B_i) \cong \mathrm{St}(\mathcal{C})$.

By Proposition~\ref{I:prop:B2}, $\hat{\mu}$ is supported on $\mathrm{pure}(\Omega)$.  Since $\mathrm{pure}$ is injective (discriminability), define the descended measure $P$ on $\sigma(\mathcal{C})$ by
\[
  P(A) \;=\; \hat{\mu}\!\bigl(\mathrm{pure}(A)\bigr)
  \qquad \text{for } A \in \sigma(\mathcal{C}).
\]
By Proposition~\ref{I:prop:B1}, $\mathrm{pure}^{-1}$ maps Baire sets in $\mathrm{St}(\mathcal{C})$ to sets in $\sigma(\mathcal{C})$; equivalently, $\mathrm{pure}$ maps $\sigma(\mathcal{C})$-sets into Baire sets, so $P$ is well-defined.  For any cylinder set $E = \mathrm{Cyl}(i, A)$,
\[
  P(E) \;=\; \hat{\mu}([E]) \;=\; \mu_i(\mathrm{Cyl}(i, A)).
\]

\textit{Uniqueness.}
Two probability measures on $(\Omega, \sigma(\mathcal{C}))$ agreeing on all cylinder sets are equal by Proposition~\ref{I:thm:obs-determination}, which requires common coarsenings.  The Stone theorem therefore assumes common coarsenings in addition to surjective evaluation, discriminability, and CE.
\end{proof}

\noindent It may appear that compactness of the Stone space eliminates the need
for CE.  Compactness yields a $\sigma$-additive Baire measure $\hat{\mu}$ on $\mathrm{St}(\mathcal{C})$, not yet on $\Omega$.  Descent to $\Omega$ is equivalent to $\hat{\mu}$ being supported on $\mathrm{pure}(\Omega)$, the principal ultrafilters; and by Proposition~\ref{I:prop:B2} this support condition is exactly CE.  In the Stone picture, CE is not an additional hypothesis but a \emph{support condition}: it is what ensures the measure concentrates on realised states rather than ideal limit points.

When the index set admits sequential common refinements and the query system discriminates points, both constructions are available simultaneously.

\begin{proposition}[Uniqueness of the observable measure]
\label{I:prop:coincidence}
Suppose the query system satisfies Surjective Evaluation, Discriminability, and sequential common refinements and common coarsenings, and let $\{\nu_i\}$ be a compatible family of probability measures.  Let $P^{\mathrm{C}}$ be the measure produced by Theorem~\ref{I:thm:obs-extension} and $P^{\mathrm{S}}$ the measure produced by Theorem~\ref{I:thm:stone-main}.  Then $P^{\mathrm{C}} = P^{\mathrm{S}}$.
\end{proposition}

\begin{proof}
Both measures agree on every cylinder set: $P^{\mathrm{C}}(\mathrm{Cyl}(i, A)) = \nu_i(A) = P^{\mathrm{S}}(\mathrm{Cyl}(i, A))$.  By observational determination (Proposition~\ref{I:thm:obs-determination}), they are equal on $\sigma(\mathcal{C})$.
\end{proof}

\begin{corollary}
The observable measure may be obtained by two distinct constructions: by countable reach in $\iota$ via sequential common refinements (Theorem~\ref{I:thm:obs-extension}), or by the Stone completion of $\mathcal{C}$ (Theorem~\ref{I:thm:stone-main}). Whenever both constructions apply, they yield the same measure.
\end{corollary}

% ------------------------------------------------------------------
\section{Synthesis}
\label{I:sec:bridge}
% ------------------------------------------------------------------

The results of this paper separate three layers.  A finitely additive charge gives a coherent valuation on the Boolean algebra of observable distinctions --- this is the finitary layer, verifiable at any single level.  Stone duality completes that algebra into the space of all coherent distinction patterns, including nonprincipal ones; compactness of the completion supplies $\sigma$-additivity there automatically.  But the completed space is larger than the realised instantiation: it contains ideal limit points, internally coherent but unrealised.  Collective exhaustion is the condition that closes the gap, ensuring the measure descends from the completion to the realised space.

Every point of $\mathrm{St}(\mathcal{C})$ is a maximal consistent assignment of yes/no answers to every observable distinction --- a coherent distinction pattern, whether or not any $\omega \in \Omega$ realises it.  The realised instantiation sits inside $\mathrm{St}(\mathcal{C})$ as the principal ultrafilters; beyond it lie ideal limit points.  In classical formulations, $\sigma$-additivity is postulated and the distinction between realised and ideal points does not arise.  In the present formulation, the distinction is visible and CE is its resolution: mass cannot persist on events that the system as a whole sees as empty.

\ifstandalone
\bibliographystyle{plain}
\bibliography{references}
\fi
